% problem-000162 4 二价铬氧化反应机理
\section*{4. 二价铬氧化反应机理}
⾦属M与单质碘在⽆⽔醋酸中反应，得到化合物 $\mathbf { M } \mathrm { I } _ { n }$ ，提纯⼲燥后测得其熔点为 $1 4 4 . 5 ^ { \circ } \mathrm { C } _ { \circ }$ 为确定 $\mathbf { M } \mathrm { I } _ { n }$ 的组成，进⾏以下实验：称取0.6263g $\mathrm { M I } _ { n } .$ ，置于特制的梨形瓶中，加⼊50 mL 6 M的盐酸，同时加⼊适量 $\mathrm { C C l } _ { 4 }$ ，⽤0.1000 M的碘酸钾溶液进⾏滴定。随着碘酸钾溶液的加⼊，可观察到 $\mathrm { C C l } _ { \downarrow }$ 层呈紫⾊，后逐渐变浅，滴⾄紫⾊褪去为终点，此时 $\mathrm { C C l _ { 4 } }$ 层呈淡⻩⾊，滴定消耗碘酸钾溶液20.00$\mathrm { m L } _ { \mathrm { c } }$ 。在 $\mathrm { C C l _ { 4 } }$ 中显⻩⾊的物种历史上曾被误认为 $\mathrm { B r } _ { 2 \mathsf { c } }$

\subsection*{4-1}
依据滴定过程中所观察到的现象，分别写出紫⾊出现（反应1）和紫⾊消失（反应2）所发⽣反应的离⼦⽅程式。写出滴定反应的总离⼦⽅程式（反应3）。

\subsection*{4-2}
通过计算，确定⾦属元素M及其碘化物 $\mathbf { M } \mathrm { I } _ { n }$ 的化学式。

\subsection*{4-3}
将 $\mathbf { M } \mathrm { I } _ { n }$ 与(CH_{3})_{3}SI按计量⽐1:2反应，得到⼀种稳定、⽆毒、⾼效的化合物T，T可⽤作染料太阳能电池中的空⽳传输材料。写出T的化学式。
